\section{Introduction}
\label{sec:introduction}

Manufacturing supervisors often lack continuous visibility into shop-floor activity.
Cameras are widely deployed, yet raw video seldom translates into timely operational understanding.
Delays between an incorrect action and corrective feedback can extend to entire shifts, which affects productivity, quality, and training effectiveness.
Industry~4.0 initiatives therefore require systems that interpret visual streams while operations unfold rather than after reports are compiled.

\visionops{} denotes an industrial visual-intelligence prototype developed under academic--industry collaboration (Alignity IQ Edge, LLC; Tec de Monterrey, Campus Estado de M{\'e}xico).
The prototype follows three operational pillars---observe, guide, and improve---and targets a \emph{digital twin of the plant} derived from multiple synchronized video sources.
Human action recognition (\har{}) constitutes a central perception module: classifying operator meta-actions such as component handling, tool use, and sheet consultation.

A recurring failure mode in multi-operator scenes is \emph{scene-level} classification, in which a single label is assigned to an entire frame although several persons perform distinct actions simultaneously.
The present study addresses this limitation through a \emph{per-person} pipeline: person detection, multi-object tracking, per-track crop buffering, frozen video embeddings, and a lightweight classifier head.
Experiments use the Industrial Human Action Recognition Dataset (\inhard{})~\cite{gaspar2020inhard}, which provides segmented RGB clips for fourteen meta-actions recorded during collaborative assembly.

The contributions of the current iteration are threefold.
First, a reproducible notebook pipeline (steps~00--06) documents exploratory data analysis, embedding extraction, classifier training, live evaluation, session logging, and quantitative model analysis.
Second, stratified sampling and configuration fingerprinting enable controlled scaling experiments (e.g., five versus one hundred clips per class) without redundant re-computation.
Third, versioned session logs capture model identity, track identifiers, confidence scores, frame snapshots, and embedding artifacts to support iteration-over-iteration comparison.

The remainder of the manuscript is organized as follows.
\Cref{sec:literature} reviews related work on industrial \har{} and video foundation models.
\Cref{sec:methods} describes the dataset, pipeline architecture, and evaluation protocol.
\Cref{sec:results} reports exploratory statistics and classifier metrics from completed runs.
\Cref{sec:discussion} interprets findings and limitations.
\Cref{sec:conclusion} summarizes conclusions and planned extensions.
